\subsection{Families of point measures}

Let X and T be two locally compact spaces, \(\pi\) a mapping of T into X, and \(g\) a finite numerical function \(\geqslant 0\) defined on T; these two functions define a mapping \(t \mapsto \lambda_t = g(t)\varepsilon_{\pi(t)}\) of T into the space \(\mathcal{M}(X)\) of measures on X, such that for every \(t \in T\) , \(\lambda_t\) is either a point measure (Ch. III, §2, No.~4) or is equal to 0. If \(f\) is a numerical function \(\geqslant 0\) defined on X, then \(\int^{*} f(x) d\lambda_t(x) = \int^{\bullet} f(x) d\lambda_t(x) = f(\pi(t))g(t)\) (recall our convention of taking this product to be 0 when \(g(t) = 0\) and \(f(\pi(t)) = +\infty\) ). Every function (with values in a topological space) defined on X is \(\lambda_t\) -measurable for every \(t \in T\) . Every mapping \(\mathbf{f}\) of X into a Banach space F is \(\lambda_t\) -integrable for all \(t \in T\) , and \(\int \mathbf{f}(x) d\lambda_t(x) = \mathbf{f}(\pi(t))g(t)\) . Finally, if \(f\) is an arbitrary numerical function defined on X, for \(f\) to be \(\lambda_t\) -integrable it is necessary and sufficient that \(f(\pi(t))g(t)\) be finite, in which case \(\int f(x) d\lambda_t(x) = f(\pi(t))g(t)\) .

DEFINITION 1. --- Let \(\mu\) be a positive measure on T. The pair \((\pi, g)\) is said to be \(\mu\) -adapted if the following conditions are satisfied:

\(1^{\circ}\) The functions \(\pi\) and \(g\) are \(\mu\) -measurable.

\(2^{\circ}\) For every function \(f \in \mathcal{K}(X)\) , the mapping \(t \mapsto f(\pi(t))g(t)\) is essentially \(\mu\) -integrable.

PROPOSITION 1. --- If the pair \((\pi, g)\) is \(\mu\) -adapted, then the mapping \(\Lambda: t \mapsto \lambda_t = g(t) \varepsilon_{\pi(t)}\) of T into \(\mathcal{M}_+(\mathrm{X})\) is scalarly essentially \(\mu\) -integrable, vaguely \(\mu\) -measurable and \(\mu\) -adequate. Conversely, if \(\Lambda\) is scalarly essentially \(\mu\) -integrable and vaguely \(\mu\) -measurable, then the function g is \(\mu\) -measurable and the restriction of \(\pi\) to the set S of \(t \in T\) such that \(g(t) \neq 0\) is \(\mu\) -measurable.

Suppose that the pair \((\pi,g)\) is \(\mu\) -adapted; for every function \(f \in \mathcal{K}(X)\) , the function \(t \mapsto \langle f, \lambda_t \rangle = f(\pi(t))g(t)\) is then essentially \(\mu\) -integrable. Let us show that \(t \mapsto \lambda_t\) is vaguely \(\mu\) -measurable. For, we first note that if \(\pi\) and g are continuous, then the mapping \(t \mapsto \lambda_t\) is vaguely continuous. In the general case, the set of compact subsets K of T such that the restrictions of \(\pi\) and g to K are continuous is \(\mu\) -dense (Ch. IV, §5, No.~10, Prop. 15); if K is such a set, then the restriction of \(t \mapsto \lambda_t\) to K is vaguely continuous, whence the first assertion of the statement. Prop. 2 of §3, No.~1 shows that \(\Lambda\) is \(\mu\) -adequate.

Conversely, suppose that \(\Lambda\) is scalarly essentially \(\mu\) -integrable and vaguely \(\mu\) -measurable; it is then \(\mu\) -adequate (§3, No.~1, Prop. 2). Since the function 1 is lower semi-continuous on X, the function \(t \mapsto \lambda_t(1) = g(t)\) is \(\mu\) -measurable (§3, Def. 1). The set S is therefore measurable (Ch. IV, §5, No.~5, Prop. 7). The set \(\mathfrak{K}\) of compact sets K ⊂ S such that \(g|K\) is continuous and \(\Lambda|K\) is vaguely continuous is \(\mu\) -dense in S (Ch. IV, §5, No.~10, Prop. 15); if K ∈ \(\mathfrak{K}\) , then the restriction to K of the mapping \(t \mapsto \varepsilon_{\pi(t)} = \frac{1}{g(t)} \lambda_t\) is therefore vaguely continuous, and this implies the continuity of \(\pi|K\) (Ch. III, §1, No.~9, Prop. 13). Since \(\mathfrak{K}\) is \(\mu\) -dense in S, the restriction of \(\pi\) to S is \(\mu\) -measurable.

We shall make use of the following lemma:

Lemma. --- Let T and X be two topological spaces, \(\pi\) a proper continuous mapping (GT, I, §10, No.~1, Def. 1) of T into X. Let g be a lower semi-continuous numerical function defined on T. For every \(x \in X\) , let \(f(x)\) be the infimum of the function \(g(t)\) in the set \(\overline{\pi}^{-1}(x)\) (infimum equal to \(+\infty\) if \(\overline{\pi}^{-1}(x) = \varnothing\) ; cf.~S, III, §1, No.~9). Then f is lower semi-continuous on X.

For every (finite) real number \(a\) , denote by \(\mathrm{B}_a\) the set of \(x \in \mathrm{X}\) such that \(f(x) \leqslant a\) , and by \(\mathrm{A}_a\) the set of \(t \in \mathrm{T}\) such that \(g(t) \leqslant a\) ; it all comes down to showing that \(\mathrm{B}_a\) is closed (GT, IV, §6, No.~2, Prop. 1). Now, \(\mathrm{A}_a\) is closed (same ref.) and the proper mapping \(\pi\) is closed (GT, I, §10, No.~1, Prop. 1); we are thus reduced to proving that \(\pi(\mathrm{A}_a) = \mathrm{B}_a\) . The obvious relation \(f(\pi(t)) \leqslant g(t)\) for all \(t \in \mathrm{T}\) implies that \(\pi(\mathrm{A}_a) \subset \mathrm{B}_a\) . On the other hand, let \(x \in \mathrm{B}_a\) ; the set \(\overline{\pi}^{-1}(x)\) is quasi-compact (GT, I, §10, No.~2, Th. 1) and nonempty, therefore there exists a \(t \in \overline{\pi}^{-1}(x)\) such that \(g(t) = \inf_{u\in \overline{\pi}^{-1} (x)}g(u) = f(x)\) (GT, IV, §6, No.~2, Th. 3); thus \(t\in \mathbf{A}_a\) and \(\pi (t) = x\).

